\documentclass[12pt,a4paper]{article}

\usepackage[utf8]{inputenc}
\usepackage[russian]{babel}
\usepackage[T2A]{fontenc}
\usepackage{amsmath, amssymb}
\usepackage{graphicx}
\usepackage{booktabs}
\usepackage{geometry}
\usepackage{hyperref}
\usepackage{natbib}
\usepackage{setspace}

\geometry{left=25mm, right=15mm, top=20mm, bottom=20mm}
\onehalfspacing
\hypersetup{colorlinks=true, linkcolor=blue, citecolor=blue, urlcolor=blue}

\title{Концентрация экономической активности в городских агломерациях:\\
эмпирический анализ на данных регионов России}
\author{И.\,И. Иванов\thanks{Институт региональной экономики; \texttt{ivanov@example.ru}}}
\date{\today}

\begin{document}
\maketitle

\begin{abstract}
\noindent В работе исследуется степень концентрации экономической активности
в крупнейших городских агломерациях России за период 2010--2023 гг.
На основе индексов Херфиндаля--Хиршмана и локализации (LQ) показано,
что рост агломерационного эффекта положительно связан с производительностью труда
($\beta = 0{,}34$, $p < 0{,}01$). Результаты имеют значение для политики
пространственного развития.

\medskip\noindent\textbf{Ключевые слова:} агломерация, региональная экономика,
индекс локализации, пространственная эконометрика.

\medskip\noindent\textbf{JEL:} R11, R12, O18.
\end{abstract}

\section{Введение}
Городские агломерации выступают драйверами экономического роста~\citep{glaeser2011}.
В России на 20 крупнейших агломераций приходится более 50\% ВВП.

\section{Данные и методология}
Используются данные Росстата по 85 субъектам РФ. Индекс локализации:
\begin{equation}
    LQ_{ij} = \frac{E_{ij}/E_i}{E_j/E},
    \label{eq:lq}
\end{equation}
где $E_{ij}$~--- занятость в отрасли $j$ региона $i$.

\section{Результаты}
Основные результаты регрессии представлены в табл.~\ref{tab:main}.

\begin{table}[htbp]
\centering
\caption{Оценки связи агломерации и производительности}
\label{tab:main}
\begin{tabular}{lcc}
\toprule
 & (1) OLS & (2) FE \\
\midrule
Индекс агломерации & $0{,}41^{***}$ & $0{,}34^{***}$ \\
                   & $(0{,}08)$     & $(0{,}09)$     \\
Плотность населения & $0{,}12^{**}$ & $0{,}10^{*}$   \\
                   & $(0{,}05)$     & $(0{,}06)$     \\
\midrule
$N$ & 1190 & 1190 \\
$R^2$ & 0{,}42 & 0{,}61 \\
\bottomrule
\end{tabular}
\end{table}

\section{Заключение}
Агломерационный эффект значим и экономически важен.

\bibliographystyle{plainnat}
\begin{thebibliography}{1}
\bibitem[Glaeser(2011)]{glaeser2011}
Glaeser E. \emph{Triumph of the City}. Penguin Press, 2011.
\end{thebibliography}

\end{document}
